\section{Dates as Constraints}
\label{sec:theory}

A year names a range of dates.
It does not identify one update date.
Our compiler keeps that range and answers questions with matching temporal scope.
It returns two evidence sets: possible support and guaranteed support.

\paragraph{Occurrence dates and replacement.}
Each claim $c$ has text and an unknown occurrence date $x_c$.
Its source bounds this date by two inclusive dates, $\ell_c\leq x_c\leq u_c$.
Dates use a common real-valued time axis.
These bounds describe uncertainty about an occurrence, rather than its duration.
The allowed assignments form the product
$\Omega=\prod_c[\ell_c,u_c]$.
A fixed text gate $G(d,c)$ accepts a direct replacement when $d$ occurs strictly after $c$.
The gate need not be transitive.
It excludes self-pairs: $G(c,c)=0$.
It must distinguish replacement from coexistence.
The compiler preserves its decisions across date assignments.

For assignment $x$, claim $c$ is active at date $t$ when
\begin{equation}
 A_x(c,t)=\mathbf{1}\!\left[
 x_c\leq t\ \land\ \nexists d:
 G(d,c),\ x_c<x_d\leq t
 \right].
\label{eq:uncertain-active}
\end{equation}
Equal dates do not establish replacement order.
A question specifies an inclusive date window $Q=[a,b]$.
Point questions set $a=b$.
A claim supports the window when it is active somewhere inside it:
\begin{equation}
 S_x(c,Q)=\mathbf{1}[\exists t\in Q:A_x(c,t)=1].
\end{equation}
Possible support requires one allowed assignment with $S_x(c,Q)=1$.
Guaranteed support requires this condition in every allowed assignment.

\paragraph{Two direct witnesses.}
One witness can force retirement despite date uncertainty.
Another can permit retirement under a particular assignment.
Their boundaries are
\begin{align}
 p_c&=\min\{u_d:G(d,c),\ \ell_d>u_c\},\label{eq:possible-boundary}\\
 r_c&=\min\{\ell_d:G(d,c),\ u_d>\ell_c\}.\label{eq:certain-boundary}
\end{align}
An empty minimum is $+\infty$.
We store one witness for each finite boundary.
These compiled certificates require a complete candidate scan.
The compiler evaluates candidate gates and updates these two minima.
It uses at most two witness records per claim.
Two records are necessary in the worst case among direct witness subsets.

\begin{theorem}[Exact window support]
\label{thm:window-support}
For every claim and inclusive query window,
\begin{align}
 \exists x\in\Omega:S_x(c,Q)=1
 &\iff \ell_c\leq b\ \land\ a<p_c,
 \label{eq:possible-window}\\
 \forall x\in\Omega:S_x(c,Q)=1
 &\iff u_c\leq b\ \land\bigl(a\leq\ell_c\ \lor\ a<r_c\bigr).
 \label{eq:guaranteed-window}
\end{align}
\end{theorem}
The first condition gives the smallest filter that never removes potentially supported evidence.
The second gives the largest filter whose retained claims always have support.
Each decision takes constant time after compilation.
Appendix~\ref{app:uncertain-proofs} gives the proofs.

\paragraph{Why window support needs a separate test.}
Suppose $x_c\in[0,10]$ and a replacing claim occurs exactly at time 5.
No date guarantees that $c$ is active.
Before time 10, $c$ might not have occurred.
Later, it might already have been replaced.
Yet $c$ supports $[0,10]$ under every assignment.
It contributes before time 5 or from its later occurrence.
Thus guaranteed window support cannot use overlap with the certainly active point set.
The order of the quantifiers matters.

\begin{corollary}[Refinement and local influence]
\label{cor:uncertain-refinement}
Narrowing date bounds can only shrink possible support and expand guaranteed support.
If two gates differ on $h$ directed pairs, either support mask changes for at most $h$ claims.
With fixed rankings, their filtered top-$k$ contexts differ in at most $2\min(k,h)$ items.
\end{corollary}
Each directed pair affects only its target claim's certificates.
When all date bounds are exact, both filters recover earliest direct replacement.
The formulas therefore extend the same policy without choosing arbitrary dates.
